\section{Integral lifting of the ramified rows}

% BEGIN A81R1
\paragraph{A81R1. Generators of integral finite cohomology.}
On a row of H80M2 let $p=\ell$, $K=E_w$, $n=[K_0:\mathbb Q_p]$,
$e=[K:K_0]=2$ and $D=ne$. Use the specified good coefficient
extension with integers $O$, uniformizer $\lambda$ and residue
field $k$. Write $T_j$ for the actual stable lattices and $V_j$
for their rational spans. The map $F:T_1\twoheadrightarrow T_2$
is equivariant and $\dim V_2=r$. Let $P_j$ be the module of
continuous $T_j$-valued cocycles whose classes become finite in
$H^1(K,V_j)$, $B_j$ the integral coboundaries, and put
$I=F(P_1)+B_2\subset P_2$. Then $P_2/B_2$ is generated over
$O$ by at most $(D+2)r$ elements.

Here are details of the integral finiteness argument. The standard
local finiteness theorem for a finite Galois module gives finite
$H^1(K,T_2/\lambda^j)$ for each $j$. Integral $H^1(K,T_2)$
is their inverse limit. Indeed representatives of compatible
classes can be chosen successively to reduce exactly to one
another by subtracting coboundaries of lifted vectors.
If a cocycle becomes a coboundary at every level, the finite
nonempty sets of primitives have nonempty inverse limit by compactness.
Thus $H^1(K,T_2)$ is compact and separated. Its quotient by
$\lambda$ injects in $H^1(K,T_2/\lambda)$ by the multiplication
sequence. Lift a basis of that finite-dimensional quotient.
Successive approximation in powers of $\lambda$ expresses every
class as an $O$-linear combination of the lifts: the errors vanish
at each fixed finite level and the coefficient series converge in $O$.
In particular integral cohomology and its submodule $P_2/B_2$ are
finitely generated. From the same exact sequence,
$H^1(K,T_2)[\lambda]$ is a quotient of $H^0(K,T_2/\lambda)$,
of $k$-dimension at most $r$. The structure theorem over $O$
therefore gives at most $r$ torsion generators of $P_2/B_2$.

The rationalization of $P_2/B_2$ is $H^1_f(K,V_2)$.
Indeed the compact image of a continuous rational cocycle lies in
$\lambda^{-j}T_2$ for some $j$, and a rational primitive similarly
acquires integral coordinates after multiplication. By the crystalline
Bloch--Kato formula of J80H1 this finite subspace has dimension
$\dim H^0(K,V_2)+\dim(D_{\rm dR}(V_2)/\mathrm{Fil}^0)$ over the
coefficient field, at most $(D+1)r$. Adding the torsion generators
proves the assertion. Choosing cocycle representatives gives the
source's generating set modulo $B_2$, with cardinality
$g\le(D+2)r\le2(D+1)r$. This argument works in every finite
local degree with the same rational dimension estimate.
% END A81R1

% BEGIN A81R2
\paragraph{A81R2. The coarse lifting test descends with no loss.}
Let $h=H$ from H80M2 and $K_1=K_0(\mu_p)$. This is a Galois
extension of $K$ of degree $s=(p-1)/e$ prime to $p$ by R79Q2.
Crystalline periods remain a basis over the extended unramified
scalars on restriction. De Rham comparison over $K_1$ extends the
filtered space by scalars, so its labelled jumps are obtained from
the restricted labels over $K$. The coefficient extension contains
the conjugates of $K_1$ by the good-layer hypothesis. Let $U$ be
the union of source jumps over its embeddings, a subset of $[-h,h]$.
For the restricted map put
\[
\begin{gathered}
 J=\big(\{a+j:a\in U,\ 1\le j\le h\}\setminus\{0\}\big)
                  \cup\{1,\ldots,h\},\\
 w(t)=t+v_p(t!),\qquad
 E=2w(4h)+\sum_{j\in J}(1+v_p(j))+w(h).
\end{gathered}
\]
On $K_1$ the cyclotomic level-one subgroup is the whole local
Galois group, and $[K_1:K_0]=p-1$. These are exactly the native
case B conditions. Its two constant height functions are $h,h$,
so its maximum lower height is $h$ and its maximum total height
is $2h$. Inserting these in the case B lifting exponent gives $E$.

Suppose the restricted integral finite-cohomology map contains
$p^E$ times every target finite class. For $c\in P_2$, restrict
its class, choose the given integral finite lift, and corestrict
this lift to $K$. Transfer commutes with $F$ and sends integral
classes to integral classes. It also preserves rational finiteness:
the map into cohomology with $B_{\rm cris}$ coefficients commutes
with transfer by its finite sum over cosets, so a zero image
remains zero. The composition of restriction and corestriction
on the original class is multiplication by $s$. As $s\in O^\times$
we recover a finite integral lift of $p^E[c]$ over $K$.
At the cocycle level the difference is in $B_2$, as required.
Thus the sole coarse lifting test here is on the cyclotomic
field itself; it is still an integral test.

Every $j\in J$ lies in $[-h,2h]\setminus\{0\}$. Summing
$1+v_p(j)$ over the entire allowed range costs $w(h)+w(2h)$.
Legendre's formula gives $v_p(t!)\le t/(p-1)\le t/2$ since
$p$ is odd. Hence $E\le2w(4h)+2w(h)+w(2h)\le18h$.
% END A81R2

% BEGIN A81R3
\paragraph{A81R3. The invariant parameter and its two primitive branches.}
For $K$ as above let $\Delta\subset\mu_{p-1}(\mathbb Z_p)$ be
the subgroup of order $s$ fixing $K$ in $K_1$, lifted by Teichm\"uller.
On each letter of L78W13 write $X$ for the chosen cyclotomic
variable, $[d](X)=(1+X)^d-1$, and $Q=\Phi_p(1+X)$. Define
\[
 Z=\frac1s\sum_{d\in\Delta}d^{-1}[d](X),\qquad v=Z^s.
\]
The substitutions satisfy $[d]\circ[c]=[dc]$, so $Z$ transforms
under $c\in\Delta$ by $c$. Its linear coefficient is one and
its constant coefficient zero. Successive inversion makes it an
integral formal coordinate. In $O[[Z]]$ a $\Delta$-fixed series
has coefficients only at multiples of $s$, since the differences
of distinct tame roots are units. Consequently $O[[X]]^\Delta=O[[v]]$
with unique coefficients, and $v=X^s$ times a unit.

Choose a primitive root $\zeta$ and put $a=v(\zeta-1)$.
Its valuation is $s/(p-1)=1/e$ by the unit just found. The element
$a$ belongs to the $\Delta$-fixed field $K$. The orbit under
$\mathrm{Gal}(K_1/K_0)$ has size $e$. Indeed for
$b\in\mathbb F_p^\times$ the quotient
$v(\zeta^b-1)/v(\zeta-1)$ reduces to $b^s$ in the residue
field, so distinct cosets modulo $\Delta$ give distinct values.
The monic polynomial $E_0(W)$ with roots these $e$ values lies in
$O_{K_0}[W]$. Its coefficients even lie in $\mathbb Z_p$, since
the series defining $v$ is defined over $\mathbb Z_p$ and the
same set of roots is stable under absolute Galois conjugation.
Its constant term has valuation one. For each coefficient letter
take $E_i$ to be the image of $E_0$. In $O[[X]]$ every
$X-(\zeta^b-1)$ divides $E_i(v(X))$ by integral division
as in W79H4. The product divides by successive evaluation and
cancellation at the remaining distinct roots. This product is $Q$.
In the residue field $E_i(v)$ has order $es=p-1$ in $X$,
because every non-leading coefficient is integral of positive
valuation, while $Q$ reduces to $X^{p-1}$. Therefore
$E_i(v)=Q$ times a unit. The Frobenius substitution
$X\mapsto XQ$ now gives $\varphi(v)=v E_i(v)^s$ times a unit.

Finally let $u=\chi(\gamma)$ for the source's level-one generator.
It satisfies $v_p(u-1)=1$. The quotient
$([u]^*v-v)/v$ is integral because $[u](X)$ is $X$ times a unit.
It vanishes at every $\zeta^b-1$ since $u\equiv1\pmod p$ and
$v(\zeta^b-1)\ne0$. It is divisible by $Q$ by the same
evaluation argument. The constant coefficient after division by
$Q$ is $(u^s-1)/p$, a unit since $p\nmid s$. Thus
$\gamma(v)-v=v E_i(v)$ times a unit. The polynomial $E_0$
has its specified root $a\in K$ with valuation $1/e$ under
every embedding. Assembling these identical letter computations
proves all clauses of $\mathrm{wachRingRelations}$ on the very
ramified rings for $e=2$ (and the proof applies to every tame
intermediate $K_0\subset K\subset K_1$). This supplies the ring
input to the scaled letter decomposition.
% END A81R3

% BEGIN A81R4
\paragraph{A81R4. Reduction of the refined exponent to the actual letters.}
Write $S,\eta$ for the integral parameters $S_0,A$ in H80M2 and
take $h'=h=H$. In particular $1\le h'$ and the positive endpoint
of every deep interval plus $S$ is at most $h'$. The strict
filtered quotient in H80M1 puts the jumps of $V_2$ in the same
deep-or-central set as those of $V_1$. Its central graded dimension
is at most $m$ from H80M2. These give the normalized hypotheses
$\mathrm{HypN}$ with the original central width and endpoints,
including the allowed degree-two alternative. At $q$ the two
embeddings extending one $K_0$-letter have the same deep sign by
R79Q1 and R79Q3. At $d$ they have opposite signs, since the
nontrivial element of $\mathrm{Gal}(K/K_0)$ interchanges the
CM type. The dual-twist reversal in H80M2 changes both signs
simultaneously and preserves this distinction.

Define $x=h/s$, $M=\lceil4x\rceil$ and
\[
\begin{gathered}
 \rho=w(2\lceil2x\rceil-1),\quad
 \sigma=w(M-1),\quad
 \kappa=e\big(w(S)+\lfloor\log_p h\rfloor\big),\\
 y_0=\lceil\eta+\tfrac12+\log_p h\rceil,
 \qquad y_2=2\sigma+\kappa+\rho+\max(2\sigma,y_0).
\end{gathered}
\]
Use $v,E_i$ from A81R3. On letter $c$ write $J_c$ for the
source's $\mathrm{JZ}$, formed inside integral target cocycles.
Explicitly, take the span of cocycles whose cohomology class is
the Herr class of a $\Delta$-invariant Wach pair $(u,z)$.
Here $u$ is supported on $c$, $E(v)^h u$ is in the Frobenius
span of $N_{c-1}$, and both $(\gamma_K-1)u$ and $z$ lie in
$v^MN$ and are $\Delta$-invariant. The pair must satisfy the
native Wach cocycle condition. This description uses the same
integral $N=N(T_2)$ and class map as in L78W13.

Suppose the following two estimates on these objects.
First, $p^{\rho+\sigma}P_2\subset B_2+\sum_cJ_c$.
Second, for each $c$ there are at most $e(e+1)m$ cocycles
$A_c\subset P_2$ such that
$p^{\sigma+y_2}J_c\subset I+\langle A_c\rangle_O$.
The latter is needed in pure position at $q$ and mixed position
at $d$, on the actual row map. Put $A_0=\bigcup_cA_c$.
There are $n$ letters, so $|A_0|\le ne(e+1)m$.
Multiplying the first inclusion by $p^{\sigma+y_2}$, applying
the second separately to each summand, and using $B_2\subset I$
gives
\[
 p^Y P_2\subset I+\langle A_0\rangle_O,\qquad
 Y=2\rho+4\sigma+\kappa+\max(2\sigma,y_0).
\]
This implication removes the quotient, normalization, divisibility
and ring-existence premises of the native reduction on these rows.
The two displayed hypotheses are the scaled finite-cocycle
decomposition and the exceptional-generator estimate at each letter.
% END A81R4

% BEGIN A81R5
\paragraph{A81R5. The two lengths and the constant 160.}
On the data of A81R1--A81R4 suppose the case B lifting test of
A81R2 and the two scaled letter estimates of A81R4 hold.
Let $Q_2=P_2/I$, $C$ the image of $\langle A_0\rangle_O$ in
$Q_2$ and $k_O=\operatorname{length}_O(O/pO)$.
The quotient $Q_2/C$ is generated by at most $g$ elements and
killed by $p^Y$. Therefore it is a quotient of $(O/p^Y O)^g$.
The submodule $C$ is generated by at most $g_0=ne(e+1)m$
elements, each killed by $p^E$ by A81R2. It is a quotient
of $(O/p^EO)^{g_0}$. Since
$\operatorname{length}_O(O/p^aO)=a k_O$ by the filtration by
powers of $p$, additivity gives the finite natural number
\[
 m_Q=\operatorname{length}_O Q_2\le(gY+g_0 E)k_O .
\]
This is the actual native $F.\mathrm{col}$ and $k_O$ is its
specified normalizer. The representatives in A81R1 allow the
quotient by $B_2$ throughout.

For the normalized rows $x=eh/(p-1)\le\tfrac32 eh/p$. From
$w(t)\le3t/2$ and $x>0$ we get
$\rho\le6x+3/2$ and $\sigma\le6x$. Also
$\kappa\le(3/2)eS+e\log_p h$ and
$y_0\le\eta+3/2+\log_p h$. Bounding the maximum by its
two nonnegative summands gives
\[
 Y\le48x+\tfrac32 eS+(e+1)\log_p h+\eta+\tfrac92
 \le80\big(eh/p+eS+e\log_p h+1+\eta\big).
\]
With $E\le18h$, $g\le2(ne+1)r$ and $g_0=ne(e+1)m$ this proves
\[
 m_Q\le160k_O\Big((ne+1)r(eh/p+eS+e\log_p h+1+\eta)
                         +ne(e+1)m(h+1)\Big).
\]
Thus the asserted explicit local premises imply the original
normalized row inequality of H80M2 with $C_A=160$, uniformly in
the good coefficient extension.
% END A81R5

% BEGIN A81R6
\paragraph{A81R6. The denominator of a rational section on a graph row.}
Suppose in addition the rational map of the row admits a
$G_K$-equivariant section $j:V_2\to V_1$. In particular this
applies to the dual of a rational direct-summand inclusion,
twisted by $1$, with the given graph lattice on the middle space.
Define $a$ to be the least nonnegative integer with
$\lambda^a j(T_2)\subset T_1$. It exists by clearing the finitely
many denominators in the images of an $O$-basis. Applying
$\lambda^a j$ to any $c\in P_2$ gives an integral cocycle in
$P_1$: the map on $B_{\rm cris}$-cohomology carries a rational
primitive to a rational primitive. Its image under $F$ is
$\lambda^a c$ as a cocycle. In particular $F.\mathrm{col}$ is
finite on each such individual row without an analytic lifting
assumption, with
\[
 \operatorname{length}_O(P_2/I)\le (ne+2)r\,a .
\]
For an integral section it is zero. The bound applies using any
known denominator $a$ and can replace the coarse test whenever
$a\le E k_O$. It also proves A81R5's final inequality directly
whenever the displayed length bound is within that budget.
Rational splitting by itself supplies no numerical comparison of
$a$ with $E k_O$: integrality of the section is measured on the
actual lattice. For the unrestricted rows of H80M2 this argument supplies
no uniform bound on section denominators. A81R7--A81R11 below
prove the scaled target decomposition independently of a section.
% END A81R6

% BEGIN A81R7
\paragraph{A81R7. A basis and the two difference equations on a ramified letter.}
Take any crystalline $T$ with jumps in $[-h,h]$ on the present
good layer. In the Taylor expansion at $1$ its integral basis
of L78W1 becomes a basis of $D_i[[Y]]$ by W79H3. The crystalline
constants are fixed by $G_K$ by W79H1. The Taylor action
$Y\mapsto dY$ thus gives the identity matrix modulo $Y$ for
every $d\in\Delta$. It gives the identity modulo $X$ on the
original basis, since $1+X=\exp(Y)$ there. Average its elements
over $\Delta$, dividing by $s\in O^\times$. The change of basis is
the identity modulo $X$, so its determinant is a unit in $O[[X]]$.
On this invariant basis $N_i^\Delta$ is free over $O[[v]]$ by A81R3.
In particular the action of the level-one generator $\gamma$ is
the identity modulo $v$ on this module: its differences are
invariant and vanish at $X=0$. On the current fields the native
$\gamma_K=\gamma$. Every cyclotomic unit occurring on $G_K$ has
$s$th power congruent to one modulo $p$, and $\chi(\gamma)^s-1$
has valuation exactly one. The minimality and choice rule for
$\gamma_K$ therefore selects $\gamma$. This proves the module
part of $\mathrm{p9}$ at $e=2$ and the action clause
$\mathrm{GamPi}$.

Put $c_0=\chi(\gamma)^s$. On $v^jN_i^\Delta/v^{j+1}N_i^\Delta$
the action of $\gamma$ is multiplication by $c_0^j$:
$\gamma(v)/v$ is invariant with constant term $c_0$, and the
basis changes by the identity modulo $v$. For $1\le j<M$
the scalar $c_0^j-1$ is a unit times $p^{1+v_p(j)}$ by the
binomial valuation formula. For $z\in vN_i^\Delta$, multiply
by this power first at $j=1$ and subtract $(\gamma-1)k_1$
with $k_1\in vN_i^\Delta$, choosing its leading coefficient to
match the residue in degree one. At each successive $j$, multiply
the residual vector and preceding correction by the new power,
and kill the coefficient in degree $j$ by an element of
$v^jN_i^\Delta$. The final discrepancy belongs to $v^MN_i^\Delta$,
and the total exponent is $\sum_{1\le j<M}(1+v_p(j))=\sigma$.
This includes $M=1$, when no correction is needed, and gives the
native $\mathrm{GradN}$.

For $M=\lceil4h/s\rceil$, W79H4 gives $Q^h\varphi N\subset N$
and $Q^hN_i\subset O[[X]]\varphi N_{i-1}$. The second inclusion
is exactly the required $\mathrm{NinLam}$ for invariant elements.
By A81R3 the first gives
$\varphi(v^MN)\subset v^M Q^{sM-h}N$.
Since $sM\ge4h$ and $h\ge1$, this is contained in
$(\lambda,X)v^MN$. Frobenius preserves $(\lambda,X)$, so its
$j$th iterate on $v^MN$ lands in $(\lambda,X)^jv^MN$.
The finite free power-series lattices are complete and separated
for that ideal. Hence for every $b\in(v^MN)^\Delta$ the convergent
sum $-\sum_{j\ge0}\varphi^j b$ uniquely solves
$(\varphi-1)z=b$ there. Indeed its partial sums telescope,
invariance is closed and preserved by Frobenius, and a fixed
vector lies in every ideal power. This proves $\mathrm{FrobPiN}$
on the given ramified module.
% END A81R7

% BEGIN A81R8
\paragraph{A81R8. A uniform power of the Wach variable lifts across a surjection.}
On any of these level-one fields, including $K_1$, fix an actual
onto map $f:T'\to T$ of crystalline lattices whose labelled
jumps lie in $[-h,h]$. Put $H_*=2h$. On every letter
\[
                 X^{H_*}N_i(T)\subseteq f N_i(T').
\]
The argument must give an exponent independent of the depth
of reduction of $O$. Write $A_j=(O/\lambda^j)[[X]]$ for $j\ge1$,
$\epsilon=1+X$ and $\varphi(X)=XQ$. This series ring is free
over its Frobenius image on $1,\epsilon,\ldots,\epsilon^{p-1}$.
Indeed modulo $\lambda$ one has $\varphi(X)=X^p$, so
$1,X,\ldots,X^{p-1}$ is a basis by collecting degrees modulo $p$.
The stated basis is the unit triangular binomial change of this
one. For surjectivity lift an expansion modulo $\lambda$ and
expand the error successively after division by $\lambda$.
For uniqueness reduce an alleged relation modulo $\lambda$ and
repeat on the coefficient lifts. Frobenius fixes $O$ and respects
these reductions. The same basis decomposition holds after
inverting $X$: $Q$ is a unit there modulo each $\lambda^j$ since
it reduces to $X^{p-1}$. Equivalently localization at $\varphi(X)$
agrees with localization at $X$, by the nilpotent coefficient ideal.

Let $L_i$ be the image of the source Wach lattice in the target
modulo $\lambda^j$ and $I_i\subset A_j$ its annihilator ideal
on the quotient. After inverting $X$ the map is onto. To verify
this first modulo $\lambda$, compare in the Witt basis of
L78W3 on the $i$th idempotent. These bases extend to invertible
comparison matrices, and $f$ between the coefficient lattices
is onto. Thus its matrix modulo $\lambda$ has full row rank
after extending from $k((X))$ to the nonzero $i$th component
of $k\otimes_{\mathbb F_p} C^\flat$. This is a faithfully flat
extension of the field, so the Laurent matrix is onto.
The field embedding and nonzero idempotent follow as in the
coefficient argument of R80D1: the reduction of the series
variable is the nonzero tilt parameter, and an idempotent
divisible by $\lambda$ in the separated Witt coefficient ring
vanishes. The same map after Laurent localization modulo
$\lambda^j$ is onto by Nakayama with the nilpotent ideal
$(\lambda)$. Consequently $X^{b_j}\in\bigcap_i I_i$ for some
$b_j$, by taking the maximum of the finitely many powers needed
to clear preimages of the target basis elements.

For $g\in I_i$ expand
$Q^{2h}g=\sum_{t=0}^{p-1}\epsilon^t\varphi(g_t)$.
Each $g_t$ belongs to $I_{i-1}$. In fact take
$z\in N_{i-1}(T)$ modulo $\lambda^j$. Then
$Q^h\varphi z\in N_i(T)$ by W79H4, and multiplication by $g$
puts this vector in $L_i$. Another $Q^h$ puts the result in
$A_j\varphi L_{i-1}$, by the other W79H4 inclusion on the source.
Write the latter element as $\sum_t\epsilon^t\varphi(z_t)$
with $z_t\in L_{i-1}$ using the series basis decomposition.
On the other hand it equals
$\sum_t\epsilon^t\varphi(g_t z)$. The components are unique:
after inverting $X$, Frobenius of the target basis is again a
basis by the two height inclusions and invertibility of $Q$.
Compare its coordinates using the Laurent decomposition.
Localization is injective on the original free module since $X$
is regular even over the truncated coefficient ring. Thus
$g_tz=z_t$ and the assertion follows.

Set $I=\bigcap_i I_i$. The cyclic permutation of the letters
shows that the coefficients of $Q^{H_*}g$ belong to $I$
whenever $g\in I$. Therefore, if
$X^{H_*}\sum_t\epsilon^t\varphi(a_t)\in I$, all
$X^{H_*}a_t$ are in $I$: multiply by $Q^{H_*}$ and use
$Q^{H_*}X^{H_*}=\varphi(X^{H_*})$. Apply this coefficient
extraction repeatedly to
$X^{H_*}(\epsilon-1)^{b_j}\in I$. Choose $d$ with $p^d>b_j$.
At each extraction the exponent in $\epsilon^k$ loses its last
base-$p$ digit since $\epsilon^p=\varphi(\epsilon)$.
After $d$ extractions one obtains $X^{H_*}$ times each
coefficient of the polynomial $(Z-1)^{b_j}$. Its leading
coefficient is one. This proves $X^{H_*}\in I$ for every $j$.
Finally the integral cokernel is finite over the noetherian
local ring $O[[X]]$ and is $\lambda$-adically separated by the
Krull intersection theorem. The stated inclusion follows by
passing to all $j$; the initial nonuniform $b_j$ does not appear.
% END A81R8

% BEGIN A81R9
\paragraph{A81R9. From the variable annihilator to a power of the local prime.}
On $K$ or $K_1$ put $a=\lceil2h/s\rceil$, with $s=1$ on
$K_1$. Assume $h\ge1$. Fix a letter $i$, write $L=N_i(T)^\Delta$
and $L'=fN_i(T')^\Delta$, meaning the image of the invariant
vectors. Since $v$ is a unit times $X^s$, A81R8 gives
$v^a L\subseteq fN_i(T')$. Averaging the preimage over
$\Delta$ gives $v^aL\subseteq L'$. On
$v^jL/v^{j+1}L$ and on $v^jL'/v^{j+1}L'$ the operator
$g=\gamma$ acts by $c_0^j$ from A81R7, for $c_0=\chi(\gamma)^s$.
For $L'$ this follows by applying $f$ to the congruence on the
source. Both filtrations are stable under $g$.

For $z\in L$ consider $v^az$ modulo $v^a L'$. It is in $L'$,
so the polynomial $P=\prod_{j=0}^{a-1}(Z-c_0^j)$ in $g$ kills it
there, successively raising the $L'$-filtration to degree $a$.
It is also in $v^a L$, so the complementary polynomial
$R=\prod_{j=a}^{2a-1}(Z-c_0^j)$ sends it into
$v^{2a}L\subseteq v^aL'$. For $0\le k<2a$ the $p$-valuation
of the Lagrange denominator at $c_0^k$ is
\[
 2a-1+v_p(k!)+v_p((2a-1-k)!)\le w(2a-1).
\]
Indeed each difference has valuation $1+v_p(j-k)$ because
$p\nmid s$. For $d=p^{w(2a-1)}$ the polynomials
$d\prod_{j\ne k}(Z-c_0^j)/\prod_{j\ne k}(c_0^k-c_0^j)$
are therefore integral. Their sum equals $d$ by interpolation.
For $k<a$ each is a multiple of $R$, and for $k\ge a$ each
is a multiple of $P$. Thus $dv^az\in v^a L'$.
Cancel the regular scalar $v^a$ to obtain $dz\in L'$.

This proves the exact $\mathrm{WKill}$ needed in A81R4 at
$d=p^\rho$ for every onto crystalline lattice map of the
specified height, including a projection of a crystalline
extension onto the trivial line. On $K_1$, $\Delta$ is trivial,
$a=2h$ and the same argument proves that $p^{w(4h-1)}$ kills
the full Wach cokernel. In particular the native cyclotomic
case B Wach exponent $w(4h)$ is available on every map entering
the cohomological lifting construction. Passing from Wach
vectors to finite integral cocycles still needs compatible
Herr pairs and their finiteness.
% END A81R9

% BEGIN A81R10
\paragraph{A81R10. Saturation at every height-one prime of the letter ring.}
Let $0\to T_0\xrightarrow{\iota}T'\xrightarrow{f}T\to0$ be an
exact sequence of full lattices in crystalline representations.
On each letter the inclusion matrix $B:N_i(T_0)\to N_i(T')$
has maximal column rank over the residue field at every height-one
prime of $A=O[[X]]$. For the prime $(\lambda)$ this follows from
injectivity on the residual lattices (the quotient is free) and the
same Witt comparison over $k((X))$ as in A81R8. This also gives full
generic rank. The rank-zero source case is immediate.

For the other primes suppose that all maximal minors have a common
irreducible factor. By Weierstrass preparation such a factor is
associated to a distinguished polynomial, and has a root $z$ in the
algebraic open unit disk. The common zero set there is finite:
any one nonzero minor, after removal of its largest power of
$\lambda$, is a unit times a distinguished polynomial.
The matrix intertwines the two integral invertible semilinear
$\gamma$-matrices. Consequently this finite set is stable under
$z\mapsto(1+z)^{\chi(\gamma)}-1$ and its inverse. A finite
orbit implies $(1+z)^{\chi(\gamma)^b-1}=1$ for some $b>0$.
The exponent is a nonzero element of $\mathbb Z_p$; after removal
of its unit part, which acts invertibly on the open unit group,
$1+z$ has $p$-power order.

At $z=0$, W79H1--W79H3 identify the specialized map with the
injective map of crystalline spaces on that letter. At
$z=\zeta-1$ of order $p$ (meaning $\zeta$ has that order), its
Taylor lattice in either representation is
$\sum_j Y^{-j}\mathrm{Fil}^jD_{\tau}[[Y]]$ by W79H2--W79H3.
Indeed these claims identify each monomial by its filtration step
and the lattice is graded. The inclusion of filtered spaces is
strict by H80M1's intersection argument applied to this sequence.
Extend an adapted basis of the subspace to a filtration-adapted
basis of the middle space. In the resulting Taylor lattices the
images of the subspace basis with the indicated $Y$ powers are
part of a basis. Thus the map has full column rank at $Y=0$.
The map at each specialization is the actual Wach map by the
morphism identity in L78W11. At a root of order $p^j$, $j>1$,
$Q(\zeta-1)\ne0$. The Frobenius matrices from the preceding
letter become invertible on evaluation there, using both height
inclusions to clear powers of $Q$; the same holds after a finite
coefficient extension containing $\zeta$. Commutation of $B$
with Frobenius transfers its rank from $\zeta^p-1$ on the
preceding letter. Induction excludes every possible common zero.
This proves the asserted rank at height one.

It follows that the image of $N_i(T_0)$ is its rational column
space intersected with $N_i(T')$. In fact solve for the preimage
by a nonzero maximal minor over $\operatorname{Frac} A$. At every
height-one localization at least one such minor is a unit, so all
coordinates are integral there. The regular local ring $O[[X]]$
is normal, hence the intersection of these localizations is $A$.
By generic ranks and $f\iota=0$ this rational column space is exactly
the kernel of $f$. The same kernel conclusion holds after applying
the Frobenius span: its columns are the Frobenius translates of
the two bases, and the inclusion matrix is $\varphi B$.
More explicitly $A$ is free on $1,\epsilon,\ldots,\epsilon^{p-1}$
over $\varphi(A)$ by lifting the expansions of A81R8 through all
coefficient powers. Hence scalar extension by $\varphi$ preserves
the exactness just proved, while the Frobenius basis maps injectively
to the actual span by the height inclusions. We obtain the exact
kernel in the Frobenius spans as well.

The images on the localized \'etale modules are still injective
on the subobject: the residual comparison in A81R8 gives an
invertible maximal minor over the residual letter field, hence a
unit there in the complete \'etale DVR. The scalars $v$ and $E(v)$
become units there since they reduce to powers of $X$ times units.
The kernel identifications therefore give the native $\mathrm{KerLam}$.
For $v^kN$, cancel $v^k$, take the preimage in $N(T_0)$ and
multiply back. For the preimage of the Frobenius span under
$E(v)^h$, take the preimage of $E(v)^h y$ in the kernel Frobenius
span, then divide by $E(v)^h$ in the \'etale module. The
idempotent-letter and $\Delta$-invariance constraints follow by
injectivity and equivariance in each case.
% END A81R10

% BEGIN A81R11
\paragraph{A81R11. The scaled letter decomposition for every finite target cocycle.}
For $c\in P_2$ use its integral extension
$T'=T_2\oplus O e_0$, where $g e_0-e_0=c(g)$.
Rational finiteness means that a $B_{\rm cris}$ primitive exists,
so subtracting it from $e_0$ extends the fixed crystalline basis
of $V_2$ by a fixed vector above one. The extension is crystalline.
Its labelled jumps are those of $V_2$ together with zero, by strict
filtered exactness, hence lie in $[-h,h]$. Normalize the Wach
generator $e$ of the actual integral trivial representation by
$1$ on its lattice, as in WLD23 and LWR5. It is fixed by Frobenius
and the cyclotomic group. In fact a Frobenius-fixed generator
compares with a constant element of $O$ by the Witt fixed-vector
identity of LWR5, and that constant must be a unit by the comparison
basis of L78W3.

By A81R9 lift $p^\rho e_i$ to $x_i\in N_i(T')^\Delta$ on each
letter, where $e_i$ here denotes the $i$th component of $e$.
The difference $(\gamma-1)x_i$ lies in $vN_i(T')^\Delta$ by A81R7
and in the kernel by its projection. A81R10 lifts it to
$z_i\in vN_i(T_2)^\Delta$. Apply A81R7 to choose
$k_i\in vN_i(T_2)^\Delta$ such that
$z_i'=p^\sigma z_i-(\gamma-1)k_i\in v^MN_i(T_2)^\Delta$.
Replace $x_i$ by $x_i'=p^\sigma x_i-\iota k_i$. It has projection
$p^{\rho+\sigma}e_i$ and $(\gamma-1)x_i'=\iota z_i'$.
Now $\varphi x_{i-1}'-x_i'$ has zero projection and is supported
at $i$. Its product with $E(v)^h$ belongs to the middle Frobenius
span by W79H4, so A81R10 gives its preimage $w_i$ with
$E(v)^hw_i$ in the kernel Frobenius span and $w_i$ invariant.
Commutation gives
\[
        (\gamma-1)w_i=\varphi z'_{i-1}-z_i'\in(v^MN(T_2))^\Delta .
\]
By A81R7 solve uniquely
$(\varphi-1)b_i=(\gamma-1)w_i$ in $(v^MN(T_2))^\Delta$.
Thus $(w_i,b_i)$ is an actual pair of the letter $J_i$ of A81R4.
Summing and using injectivity of $\varphi-1$ on the deep module
gives $\sum_i b_i=\sum_i z'_i$.

Write $x'=\sum_i x_i'$, $d=p^{\rho+\sigma}$ and use the original
Witt comparison embedding $j_D$. The Witt vector
$b=j_D(x')-d(1\otimes e_0)$ lies in the Witt tensor with $T_2$
by the integral splitting of the underlying lattices. It satisfies
$(\varphi-1)b=j_D(\sum_i w_i)$ after identifying the kernel.
For an element $g$ whose cyclotomic image is trivial or tame,
$g$ fixes $x'$, so $-\delta b(g)=d(1\otimes c(g))$.
For a lift $g$ of the chosen procyclic generator,
$j_D(\sum_i b_i)-\delta b(g)=d(1\otimes c(g))$, by the
sum identity and $(\gamma-1)x'=\iota\sum_i z_i'$.
These are precisely the two rules determining LWR5's continuous
Herr cocycle. By its additivity the sum of the classes of
$(w_i,b_i)$ is $d[c]$. Choosing integral cocycle representatives
of the separate classes therefore proves
\[
              p^{\rho+\sigma}P_2\subseteq B_2+\sum_iJ_i .
\]
This supplies the first scaled estimate of A81R4 for the original
ramified rows, including the mixed prime. In the proof of A81R5
the remaining per-letter estimate concerns lifts through the
particular row map with at most $6m$ exceptions per letter.
% END A81R11
